\chapter{Annales — Sujets complets type BEPC}

\noindent\textit{Ce chapitre propose trois épreuves complètes de Sciences de la Vie et
de la Terre, de l'Éducation à l'Environnement et à la Santé, construites pour
reproduire fidèlement la structure et le barème réels de l'épreuve du BEPC (Partie A —
Évaluation des ressources : I. Connaissances et II. Compétences méthodologiques ;
Partie B — Évaluation des compétences), telle que décrite au chapitre « Méthodologie
de l'épreuve ». Contrairement aux annales officielles d'une session précise, ces trois
sujets sont des créations originales du fascicule, chacun croisant plusieurs chapitres
du programme autour d'une situation concrète, avec un corrigé entièrement rédigé et
vérifié.}
\vspace{8pt}

\connecteBanner

% ═══════════════════════════════════════════════════════════════
\section{Sujet 1 — Une campagne de santé au village}

\noindent\textbf{Examen :} BEPC \quad\textbf{Épreuve :} SVTEEHB \quad\textbf{Notée
sur :} 20 points.
\par\smallskip\noindent
L'épreuve comporte deux parties A et B, indépendantes et obligatoires.

\subsection*{Partie A — Évaluation des ressources (10 points)}

\paragraph{I. Connaissances (4 points)}
\begin{enumerate}[label=\arabic*)]
\item Définir la vaccination. \hfill (1pt)
\item Citer les quatre signes caractéristiques de la réaction inflammatoire. \hfill
(1pt)
\item Qu'est-ce qu'un allergène ? \hfill (1pt)
\item Citer les trois voies réelles de transmission du VIH. \hfill (1pt)
\end{enumerate}

\paragraph{II. Compétences méthodologiques (6 points)}
\noindent\textbf{Document :} Un élève reçoit une première injection d'un vaccin (jour
J0), puis une seconde injection du même vaccin, ou rappel (jour J30). La quantité
d'anticorps spécifiques présents dans son sang évolue ainsi : quasiment nulle avant J0
; faible et d'apparition lente après la première injection ; élevée et d'apparition
rapide après le rappel de J30.
\begin{enumerate}[label=\arabic*)]
\item Que représente l'apparition d'anticorps après la première injection (J0) ?
\hfill (2pt)
\item Comparer la vitesse et l'intensité de la réponse observée après le rappel (J30)
à celles observées après la première injection. \hfill (2pt)
\item Expliquer cette différence à l'aide de la notion étudiée en cours. \hfill (2pt)
\end{enumerate}

\subsection*{Partie B — Évaluation des compétences (10 points)}

\noindent\textbf{Situation :} Lors d'une campagne de santé dans un village, une
infirmière vaccine les élèves contre le tétanos, dépiste des cas de paludisme, et
prend en charge un élève victime d'une réaction allergique après une piqûre
d'insecte.

\begin{enumerate}[label=\textbf{Tâche \arabic*.}, leftmargin=2.2cm]
\item Expliquer pourquoi la vaccination contre le tétanos assure une protection
durable, alors qu'un sérum antitétanique agirait immédiatement mais pour une durée
plus courte. \hfill (3pt)
\item Un élève présente fièvre et frissons après avoir été plusieurs fois piqué par
des moustiques. Quelle maladie l'infirmière doit-elle suspecter en priorité ? Quel
geste de prévention aurait pu limiter ce risque ? \hfill (3pt)
\item Un autre élève, déjà piqué par une guêpe un an plus tôt sans réaction
particulière, développe cette fois une éruption cutanée et des démangeaisons après
une nouvelle piqûre. Ce phénomène est-il compatible avec une réaction allergique ?
Justifier en s'appuyant sur le mécanisme en deux temps. \hfill (2pt)
\item Proposer, en une phrase chacun, un conseil de prévention adapté à chacune des
trois situations rencontrées par l'infirmière (vaccination, paludisme, allergie).
\hfill (2pt)
\end{enumerate}

% ═══════════════════════════════════════════════════════════════
\section{Sujet 1 — Corrigé}

\subsection*{Partie A}

\paragraph{I. Connaissances}
\begin{enumerate}[label=\arabic*)]
\item La vaccination consiste à injecter une forme inactivée ou atténuée d'un microbe
ou de l'une de ses toxines, incapable de provoquer la maladie mais capable de
déclencher la réponse spécifique de l'organisme.
\item Rougeur, chaleur, douleur, gonflement.
\item Un allergène est une substance normalement inoffensive qui déclenche, chez
certaines personnes, une réaction immunitaire excessive.
\item Voie sexuelle, voie sanguine, voie mère-enfant.
\end{enumerate}

\paragraph{II. Compétences méthodologiques}
\begin{enumerate}[label=\arabic*)]
\item Cela représente la mise en place d'une réponse immunitaire spécifique contre le
microbe (ou la toxine) contenu dans le vaccin, avec fabrication d'anticorps et de
cellules mémoires.
\item Après le rappel, la réponse est beaucoup plus \textbf{rapide} et beaucoup plus
\textbf{intense} qu'après la première injection.
\item Cette différence s'explique par la \textbf{mémoire immunitaire} : lors de la
première injection, des cellules mémoires se sont mises en place ; lors du rappel,
ces cellules déjà présentes permettent une réponse spécifique immédiate et efficace.
\end{enumerate}

\subsection*{Partie B}

\textbf{Tâche 1.} La sérothérapie apporte des anticorps déjà fabriqués, qui agissent
immédiatement mais finissent par disparaître de l'organisme (protection de courte
durée). La vaccination stimule l'organisme à fabriquer lui-même ses anticorps et des
cellules mémoires durables : la protection met plus de temps à s'installer, mais dure
beaucoup plus longtemps.

\textbf{Tâche 2.} Il faut suspecter un \textbf{paludisme} (fièvre et frissons après
des piqûres de moustique). Dormir sous une moustiquaire imprégnée (ou éliminer les
eaux stagnantes autour de l'habitation) aurait pu limiter ce risque.

\textbf{Tâche 3.} Oui, c'est compatible : la première piqûre a provoqué une
\textbf{sensibilisation} silencieuse, sans symptôme visible ; lors de la seconde
piqûre, l'organisme déjà sensibilisé libère massivement de l'histamine, provoquant
les symptômes observés.

\textbf{Tâche 4.} Vaccination : respecter le calendrier vaccinal et faire les
rappels recommandés. Paludisme : dormir sous une moustiquaire imprégnée et éliminer
les eaux stagnantes. Allergie : éviter tout contact avec l'allergène identifié.

% ═══════════════════════════════════════════════════════════════
\section{Sujet 2 — Le carnet de santé d'un nouveau-né}

\noindent\textbf{Examen :} BEPC \quad\textbf{Épreuve :} SVTEEHB \quad\textbf{Notée
sur :} 20 points.

\subsection*{Partie A — Évaluation des ressources (10 points)}

\paragraph{I. Connaissances (4 points)}
\begin{enumerate}[label=\arabic*)]
\item Définir un gène. \hfill (1pt)
\item Combien de chromosomes compte un caryotype humain normal ? \hfill (1pt)
\item Qu'est-ce que la mitose ? \hfill (1pt)
\item Citer les trois génotypes possibles pour la drépanocytose. \hfill (1pt)
\end{enumerate}

\paragraph{II. Compétences méthodologiques (6 points)}
\noindent\textbf{Document :} Le caryotype d'un nouveau-né, établi à partir d'une
prise de sang, montre 44 autosomes et deux chromosomes sexuels X et Y. Ses parents
sont tous deux porteurs sains de la drépanocytose (génotype AS).
\begin{enumerate}[label=\arabic*)]
\item a) De quel sexe s'agit-il ? b) Combien de chromosomes ce caryotype compte-t-il
au total ? c) Ce caryotype présente-t-il une anomalie chromosomique ? Justifier.
\hfill (3pt)
\item a) Construire le tableau de croisement du couple AS $\times$ AS. b) Quelle est
la probabilité que cet enfant soit atteint de drépanocytose (SS) ? \hfill (3pt)
\end{enumerate}

\subsection*{Partie B — Évaluation des compétences (10 points)}

\noindent\textbf{Situation :} Un couple attend un enfant. Le père est porteur sain de
la drépanocytose (génotype AS) et de groupe sanguin O. La mère ne connaît pas son
statut pour la drépanocytose ; elle est de groupe sanguin A.

\begin{enumerate}[label=\textbf{Tâche \arabic*.}, leftmargin=2.2cm]
\item Rappeler ce que signifie, pour le père, le génotype AS. \hfill (1pt)
\item Si la mère est de génotype AA (non porteuse), l'enfant peut-il être atteint de
drépanocytose (SS) ? Construire le tableau de croisement correspondant pour
justifier. \hfill (3pt)
\item Le père est de génotype OO (groupe O) et la mère de génotype AO (groupe A).
Construire le tableau de croisement du système ABO et donner les groupes sanguins
possibles pour l'enfant. \hfill (3pt)
\item Expliquer pourquoi un suivi du caryotype pendant la grossesse permet de
détecter une éventuelle trisomie 21 avant la naissance. \hfill (3pt)
\end{enumerate}

% ═══════════════════════════════════════════════════════════════
\section{Sujet 2 — Corrigé}

\subsection*{Partie A}

\paragraph{I. Connaissances}
\begin{enumerate}[label=\arabic*)]
\item Un gène est un segment d'ADN, situé à un locus précis d'un chromosome, qui
porte l'information nécessaire à un caractère héréditaire donné.
\item 46 chromosomes.
\item La mitose est le mode de division d'une cellule qui produit deux cellules
filles génétiquement identiques entre elles et identiques à la cellule mère.
\item AA (sain), AS (porteur sain), SS (malade).
\end{enumerate}

\paragraph{II. Compétences méthodologiques}
\begin{enumerate}[label=\arabic*)]
\item a) C'est un garçon (chromosomes XY). b) 44 $+$ 2 $=$ 46 chromosomes. c) Non,
aucune anomalie : 46 chromosomes avec une paire sexuelle XY correspond à un caryotype
masculin normal.
\item a) Le tableau de croisement AS $\times$ AS donne les quatre combinaisons AA,
AS, AS, SS. b) La probabilité que l'enfant soit SS (atteint) est de $\dfrac14$, soit
$25\,\%$.
\end{enumerate}

\subsection*{Partie B}

\textbf{Tâche 1.} Le génotype AS signifie que le père possède un exemplaire normal
(A) et un exemplaire muté (S) du gène de l'hémoglobine : il est porteur sain, sans
symptôme de la maladie.

\textbf{Tâche 2.} Non : le père AS transmet soit A soit S ; la mère AA ne peut
transmettre que A. Le tableau de croisement donne uniquement des enfants AA ou AS,
jamais SS. L'enfant ne peut donc pas être atteint de drépanocytose dans ce cas.

\textbf{Tâche 3.} Le père OO ne transmet que l'allèle O. La mère AO transmet soit A
soit O. Le tableau de croisement donne donc AO (groupe A) ou OO (groupe O) : l'enfant
peut être de groupe A ou de groupe O.

\textbf{Tâche 4.} Le caryotype établi à partir d'une analyse prénatale (prélèvement
de liquide amniotique) permet de compter les chromosomes du fœtus ; la présence de
trois chromosomes 21 au lieu de deux (47 chromosomes au total) signale une trisomie
21, détectable avant la naissance.

% ═══════════════════════════════════════════════════════════════
\section{Sujet 3 — Une forêt menacée et une rivière polluée}

\noindent\textbf{Examen :} BEPC \quad\textbf{Épreuve :} SVTEEHB \quad\textbf{Notée
sur :} 20 points.

\subsection*{Partie A — Évaluation des ressources (10 points)}

\paragraph{I. Connaissances (4 points)}
\begin{enumerate}[label=\arabic*)]
\item Qu'est-ce qu'un écosystème ? \hfill (1pt)
\item Citer, dans l'ordre, les quatre étapes de formation d'une roche sédimentaire.
\hfill (1pt)
\item Quel est l'agent responsable du paludisme et comment se transmet-il ? \hfill
(1pt)
\item Citer deux menaces humaines qui pèsent sur les forêts. \hfill (1pt)
\end{enumerate}

\paragraph{II. Compétences méthodologiques (6 points)}
\noindent\textbf{Document :} Lors de travaux près d'une rivière, trois strates ont
été mises à jour. La strate 1 (la plus basse) contient des fossiles d'organismes
marins. La strate 2 (intermédiaire) contient des fossiles d'organismes terrestres. La
strate 3 (la plus haute) ne contient aucun fossile.
\begin{enumerate}[label=\arabic*)]
\item Classer ces trois strates de la plus ancienne à la plus récente, en justifiant
à l'aide d'un principe étudié en cours. \hfill (2pt)
\item Que peut-on déduire du milieu qui existait à cet endroit au moment du dépôt de
la strate 1 ? Justifier. \hfill (2pt)
\item Une analyse de l'eau de cette rivière révèle la présence de bactéries
pathogènes. Quelle voie de contamination cela représente-t-il pour les villageois qui
consomment cette eau sans la traiter ? \hfill (2pt)
\end{enumerate}

\subsection*{Partie B — Évaluation des compétences (10 points)}

\noindent\textbf{Situation :} Une ONG étudie une forêt tropicale traversée par une
rivière. Elle constate une déforestation croissante due à l'exploitation du bois, une
rivière polluée par des bactéries, et plusieurs cas de paludisme chez les villageois
vivant près de zones d'eaux stagnantes laissées par la déforestation.

\begin{enumerate}[label=\textbf{Tâche \arabic*.}, leftmargin=2.2cm]
\item Expliquer pourquoi la déforestation peut favoriser la prolifération des
moustiques responsables du paludisme. \hfill (3pt)
\item Proposer deux mesures pour protéger cette forêt et limiter la déforestation.
\hfill (2pt)
\item Expliquer pourquoi consommer l'eau de cette rivière sans la traiter présente un
risque de contamination, et préciser la voie concernée. \hfill (3pt)
\item Expliquer en quoi la disparition progressive de cette forêt pourrait affecter
la biodiversité locale. \hfill (2pt)
\end{enumerate}

% ═══════════════════════════════════════════════════════════════
\section{Sujet 3 — Corrigé}

\subsection*{Partie A}

\paragraph{I. Connaissances}
\begin{enumerate}[label=\arabic*)]
\item Un écosystème est l'ensemble formé par un milieu et les êtres vivants qui y
habitent, en interaction les uns avec les autres et avec leur milieu.
\item Altération, transport, sédimentation, diagenèse.
\item Un parasite du genre \textit{Plasmodium}, transmis par la piqûre d'un
moustique anophèle femelle.
\item Par exemple : la déforestation et le braconnage.
\end{enumerate}

\paragraph{II. Compétences méthodologiques}
\begin{enumerate}[label=\arabic*)]
\item D'après le principe de superposition, la strate la plus basse est la plus
ancienne. Ordre de la plus ancienne à la plus récente : strate 1, puis strate 2, puis
strate 3.
\item La présence de fossiles marins dans la strate 1 indique que cet endroit était
autrefois recouvert par la \textbf{mer} au moment du dépôt de cette strate.
\item Cela représente une contamination par \textbf{voie digestive} (ingestion d'eau
contaminée par des bactéries pathogènes).
\end{enumerate}

\subsection*{Partie B}

\textbf{Tâche 1.} La déforestation laisse le sol nu et modifie l'écoulement de
l'eau, ce qui favorise la formation d'eaux stagnantes ; ces eaux stagnantes
constituent le lieu de reproduction des moustiques anophèles, favorisant ainsi la
prolifération des moustiques et donc la transmission du paludisme.

\textbf{Tâche 2.} Par exemple : créer une réserve naturelle protégée, et adopter des
lois limitant l'exploitation forestière (ou pratiquer le reboisement).

\textbf{Tâche 3.} L'eau contaminée par des bactéries pathogènes, si elle est bue sans
traitement, fait pénétrer ces bactéries dans l'organisme par \textbf{voie digestive}
: c'est une contamination directe, pouvant provoquer une infection (par exemple des
diarrhées).

\textbf{Tâche 4.} La disparition de la forêt entraîne la perte d'habitats et de
sources de nourriture pour de nombreuses espèces, qui dépendent les unes des autres
(chaînes alimentaires, pollinisation, décomposeurs) : la disparition d'une espèce ou
d'un habitat peut donc déséquilibrer tout le réseau d'interdépendances de
l'écosystème forestier.
